\section{관련 연구}

\paragraph{스크린샷 기반 GUI 이해와 그라운딩.}
스크린샷 기반 GUI 모델은 구조화된 DOM 없이 화면에서 요소를
찾고 그 의미를 이해하도록 학습한다. SeeClick은 GUI grounding
continual pre-training 뒤에 downstream agent를 적응시키며,
중간 grounding 향상과 agent 성능을 함께 추적했다
\citep{cheng2024seeclick}. GUICourse는 웹 화면의 텍스트, 좌표,
레이아웃 순서를 장문 target으로 예측하고 데이터 규모에 따른 OCR,
grounding, navigation 향상을 보고했으며
\citep{chen2025guicourse}, MultiUI는 webpage caption을 포함한
여러 이해 과제를 함께 학습해 웹과 문서 이해로의 전이를 보였다
\citep{liu2025multiui}. 이 결과들은 화면 supervision의 유용성을
보이지만, 같은 화면의 information budget을 단계적으로 비교하지
않아 description coverage의 효과가 데이터 규모와 과제 mixture에
섞여 있다.

기존 GUI caption 연구는 상세함이 반드시 긴 문장을 뜻하지 않음을
보여준다. AutoGUI는 전후 화면에서 만든 기능 설명을 element HTML과
축약 설명에 비교해 기능ㆍaffordance 정보의 이점을 분석했고
\citep{li2025autogui}, DeskVision은 평균 12.2자의 region caption에
요소 type, visible text, attribute를 담아 grounding을 개선했다
\citep{xu2025deskvision}. 본 연구는 이들과 달리 whole-page
description을 대상으로 하며, 동일한 구조화 사실에서 보존되는
정보의 범위를 통제한다.

\paragraph{GUI agent의 단계적 학습.}
GUI agent 학습은 지각ㆍ그라운딩과 행동ㆍtrajectory supervision을
단계적으로 연결하거나 하나의 mixture에서 결합한다. UI-TARS는
element description과 dense caption을 상호작용 궤적과 함께
사용하고, MolmoWeb은 대규모 screenshot QA와 grounding을
trajectory 중심 SFT mixture에 포함한다
\citep{qin2025uitars,gupta2026molmoweb}. Aguvis와 InfiGUIAgent는
grounding 중심 단계와 planning 또는 expectation--reflection
단계를 구분하여 중간 능력과 최종 수행을 분석한다
\citep{xu2025aguvis,liu2026infiguiagent}. 이들 연구는 단계적 평가의
필요성을 보여주지만 초기 description의 상세도를 독립 변수로
다루지는 않는다.

후속 trajectory의 품질도 최종 점수를 크게 바꾼다. WebSTAR는
성공한 trajectory 안의 step을 다시 채점해, 97K 전체 step보다
선별한 46K step을 학습한 모델이 더 높은 성공률을 보임을
보고했다 \citep{he2026webstar}. 또한 실제 전후 화면으로
forward/inverse dynamics를 학습하는 방법은 상태 변화 supervision의
효용을 보이지만 \citep{liu2026statetransition}, 본 연구의 현재
화면에 대한 텍스트 기반 예상 결과와는 다르다. 이에 우리는 모든
조건에 동일한 IT, Action, Trajectory sample과 loss mask를 적용하고,
각 단계의 변화량을 보고하여 초기 description의 효과와 후속
학습의 효과를 분리한다.
